\chapter{绪论}

\section{研究背景与意义}

随着云计算、大数据和人工智能技术的广泛应用，计算资源逐渐从本地计算环境转移到云端数据中心。云计算技术通过资源池化和按需分配机制显著提高了计算效率和成本效益，因此被各大云提供商广泛采用。然而，数据控制权与物理计算环境的分离带来了新的安全挑战。由于用户在将数据和应用迁移到云端后不再完全掌控底层硬件和系统软件环境，这导致传统基于边界防护的安全模型难以适应当前的计算范式\cite{lee2013view}。

传统信息安全技术主要侧重于“传输中数据”和“存储中数据”的安全性。例如，TLS 协议和磁盘加密技术能够有效确保数据在网络传输与持久化存储过程中的机密性。然而，对于“使用中数据”的保护，一直以来缺乏有效的解决方案。在当前云计算的虚拟化与多租户环境中，机密数据在内存中的处理过程可能暴露于高权限软件，甚至物理攻击者的威胁之下\cite{win2014virtual}。

在此背景下，机密计算作为一种新兴的安全计算范式开始崭露头角。机密计算通过利用硬件支持的可信执行环境（Trusted Execution Environment, TEE），为应用程序提供隔离的运行空间，从而在一定程度上减轻了对操作系统或虚拟机监控器的信任依赖\cite{sabt2015trusted}，有效保障了“使用中数据的安全”。近年来，机密计算被视为确保“数据在使用中”安全的关键技术方向，已在金融风控、隐私计算、跨机构数据协作以及人工智能模型训练等场景中受到广泛关注\cite{confidentialcomputingconsortium2022}。随着产业界和学术界的持续投入，机密计算逐步从概念验证阶段走向规模部署\cite{agarwal2025confidential}。在机密计算的技术体系中，虚拟化场景下的安全问题尤为突出。云计算平台广泛使用虚拟机技术来实现资源隔离和多租户部署，但传统虚拟化模型默认虚拟机监控器（Hypervisor）处于完全可信的状态\cite{garfinkel2003virtual}。一旦虚拟机管理程序被攻破或内部人员滥用权限，虚拟机内的敏感数据可能面临泄露的风险。

为解决该安全问题，AMD 提出了 SEV（Secure Encrypted Virtualization）技术，该 TEE 技术提供硬件支持的机密虚拟机保护机制。AMD SEV 提供处理器内存加密能力，通过在内存控制器中集成加密引擎，为每个虚拟机分配独立的加密密钥，使虚拟机内存以加密形式存储于物理内存中\cite{kaplan2016amd}。这种内存加密确保即使虚拟机管理程序或底层系统软件具有最高权限，也无法直接读取虚拟机的明文内存数据，从而在体系结构层面强化了虚拟化环境的机密性。

然而，已有研究表明，机密计算体系不具备绝对安全性。虽然硬件隔离和内存加密机制理论上提供了更强的保护边界，系统安全性仍然取决于复杂的微架构实现和共享资源管理。近年来，许多研究揭示了针对可信执行环境的攻击路径，如侧信道攻击、瞬态执行攻击以及由资源共享引发的信息泄露问题\cite{bulck2018foreshadow}\cite{lou2022survey}。这些安全研究表明，即使在硬件增强的安全模型下，微架构层面的共享特性仍可能引发新的安全风险。现代多核处理器通常采用复杂的缓存层级结构和一致性协议，以维持数据一致性。这些机制在提高系统效率的同时，也带来了新的可观察行为特征。已有研究对微架构时序攻击和共享资源攻击进行了系统总结，指出缓存、分支预测和一致性机制可能成为攻击者推断敏感信息的媒介\cite{ge2018survey}。

因此，围绕 AMD SEV 平台在多核共享环境下的底层微架构交互行为开展研究，有助于更清晰地认识机密虚拟化体系的潜在风险。现代多核处理器复杂的缓存层级结构、一致性维护协议以及内存总线仲裁机制，在提高系统效率的同时，也带来了新的可观察行为特征。

基于此，本文将聚焦于 AMD SEV 体系下的缓存一致性机制与内存竞争所引发的缓存侧信道问题，从体系结构与系统安全角度分析其机理与攻击能力，对机密计算环境中的新型微架构侧信道风险进行系统研究，为相关防御机制的改进提供参考。

\section{国内外研究现状}

伴随机密计算技术的快速发展，其安全性问题也逐渐成为学术界与工业界关注的重点。国内外研究者围绕可信执行环境与机密虚拟化平台展开了大量安全分析工作，在攻击和防御领域都取得了不少成就。本节将从机密计算平台的总体安全研究出发，综述当前针对机密计算的瞬态执行攻击、传统缓存侧信道攻击等微架构攻击研究，并进一步分析近年来逐渐受到关注的基于缓存一致性（Cache Coherence）的侧信道攻击研究进展，最后总结现有工作的不足并引出本文的研究定位。

\subsection{机密计算安全研究进展}
机密计算通过硬件机制构建可信执行环境，以在不完全可信的系统软件环境中保护运行时数据安全。目前主流的机密计算实现包括 Intel SGX、AMD SEV 以及 ARM CCA 等\cite{kaplan2016amd,li2022design,costan2016intel}。这些技术通过内存加密、访问控制以及硬件隔离等安全机制限制操作系统或虚拟机监控器对敏感数据的访问，从而实现对应用执行环境的保护。

早期对于机密计算平台的安全研究主要集中于逻辑层攻击。在早期机密计算虚拟机平台中，系统缺乏完整性保护机制，攻击者可以通过操控页表映射实现对虚拟机内存内容的读取。Morbitzer等提出的 SEVered 攻击通过重新映射虚拟机内存页，使攻击者能够逐步恢复受保护虚拟机中的数据内容\cite{morbitzer2018severed}。这些研究指出了早期机密虚拟化技术在内存完整性保护方面的不足，并推动了后续 SEV-ES 与 SEV-SNP 等安全机制的发展。

然而，即便在引入了更强完整性保护的 SNP 模型下，新的安全问题依然层出不穷。例如，Li 等人针对早期的设计假设，提出了 CrossLine 攻击打破了机密虚拟机基于“崩溃即安全”（Security-by-Crash）的内存隔离\cite{li2021crossline}。针对 SEV-SNP 的具体实现与硬件特征，Bagia 等人分析指出，SEV-SNP 中的 RMP（反向映射表）条目缓存机制会产生微架构层面的状态残留，这为新型信息泄露提供了潜在途径\cite{bagia2025close}；Schl{\"u}ter 等人深入挖掘了处于底层特权协议逻辑与微架构交互中的缺陷，接连提出了 RMPocalypse\cite{schlüter2025rmpocalypse} 以及 Heracles\cite{schlüter2025heracles} 选择明文攻击。不仅如此，针对复杂系统机制交互的漏洞挖掘也取得了诸多进展。例如，Schl{\"u}ter 等人还在 WeSee\cite{schluter2024wesee} 与 Heckler\cite{schluter2024heckler} 工作中，证明了恶意宿主机可以通过精心构造的 \#VC 异常和恶意中断注入突破机密虚拟机的屏障；Zhang 等人提出的 CacheWarp\cite{zhang2024cachewarp} 揭示了利用状态回滚实现的软件形式故障注入漏洞；而 Gast 等人的 CounterSEVeillance\cite{gast2025counterseveillance} 则演示了如何利用性能计数器实现细粒度侧信息窃取。这些工作充分表明，机密计算在应对拥有底层系统控制权的高权限敌手时，其设计与实现仍面临极大挑战。

随着机密计算技术的不断演进，研究者逐渐将关注点从系统软件层面的攻击扩展到处理器微架构层面。研究表明，即使在硬件安全隔离机制存在的情况下，处理器内部的共享资源仍可能形成潜在的信息泄露通道 \cite{ge2018survey,lou2022survey}。例如，多核处理器中的缓存结构、分支预测器以及地址转换缓存等组件通常被多个处理器核心共享，其状态变化可能通过执行时间差异或访问行为被攻击者观测，从而形成侧信道攻击。因此，如何评估和缓解微架构层面的信息泄露问题，逐渐成为机密计算安全研究的重要方向。

\subsection{微架构侧信道攻击研究}
微架构侧信道攻击是当前硬件安全研究中最前沿的领域之一。此类攻击不直接读取受保护数据，而是通过分析处理器内部共享资源的访问行为推断敏感信息。目前的主要研究方向上是瞬态执行攻击和缓存侧信道攻击。

瞬态执行攻击（Transient Execution Attacks）是近年来最具影响力的一类微架构攻击。Spectre 与 Meltdown 系列攻击揭示了现代处理器在推测执行与乱序执行过程中可能访问本不应暴露的数据，并在架构状态回滚后留下微架构痕迹，例如缓存状态变化 \cite{kocher2020spectre,lipp2018meltdown}。攻击者可以通过观测这些变化从而恢复敏感数据。Van Bulck 等提出的 Foreshadow 攻击进一步证明，即使在可信执行环境中，利用瞬态执行漏洞仍然可以突破 SGX enclave 的隔离保护并读取其中的敏感数据 \cite{bulck2018foreshadow}。

除瞬态执行攻击之外，缓存侧信道攻击也是机密计算研究中的重要方向。Osvik 等提出的 Prime+Probe 攻击通过构造缓存集合冲突并测量访问延迟，实现对目标程序缓存访问模式的推断 \cite{osvik2006cache}。Yarom 等提出的 Flush+Reload 攻击利用共享内存页实现更高分辨率的信息泄露，该方法能够精确检测目标程序对特定缓存行的访问行为，并已被用于恢复加密算法密钥 \cite{yarom2014flushreload}。而后续研究进一步证明，在云计算环境中即使虚拟机之间不存在直接通信，也可以通过共享最后一级缓存（LLC）构造跨虚拟机侧信道攻击 \cite{ristenpart2009hey}。

此外，一些研究还探索了其他共享硬件资源带来的侧信道，例如分支预测器攻击、TLB 侧信道以及执行单元竞争攻击等 \cite{gruss2016flushflush,lou2022survey}。这些研究表明，现代处理器为提高性能而设计的共享结构往往具有可观测行为，而此类可观测的行为可能会被攻击者利用，从而构造隐蔽的信息传输通道。

在机密计算的特殊威胁模型下，不可信实体（如恶意宿主机）拥有的系统调度权限使得微架构攻击的威力被进一步放大。例如，Wilke 等人设计的 SEV-Step 框架通过巧妙利用计时器中断与异常机制，使得恶意 Hypervisor 能够以指令级粒度对机密虚拟机执行单步攻击（Single-Stepping）\cite{wilke2024sevstep}，为高精度侧信道测量提供了理想的观测基础。此外，传统防御观念往往假定不可完全解密的硬件加密机制足以隔绝敏感数据，然而 Li 等人开展的系统性分析表明，即便是完全密态的数据状态，在结合高分辨率的宿主机侧信道测量后，依然会产生严重的密文侧信道（Ciphertext Side Channels）信息泄露\cite{li2022systematic}。这类针对特定机密计算架构环境的探索，揭示了处于特权级下的微架构侧信道突破数据机密性保证的严重性。

\subsection{缓存一致性侧信道研究}
在多核处理器体系结构中，缓存一致性协议（Cache Coherence Protocol）用于维护不同处理核心之间的数据一致性。常见的一致性协议包括 MESI 和 MOESI 等，其核心机制是通过缓存行状态转换和失效消息传播来保证不同核心看到一致的数据视图 \cite{amd64manual}。

在机密计算环境中，缓存一致性相关行为同样可能成为新的信息泄露来源。随着 SEV 等机密虚拟化技术的广泛部署，研究者开始关注在机密虚拟机环境中重新构造高分辨率缓存侧信道的可能性。由于机密计算平台通常限制共享内存以及部分缓存操作指令，传统 Flush+Reload 等攻击在机密虚拟机环境中难以直接实现，因此研究者开始探索新的微架构信号来源。

Chiang等在 Reload+Reload 工作中，通过对 AMD SEV 处理器微架构行为进行分析，发现了两种此前未被公开的侧信道机制：缓存刷新侧信道（RRFS）和内存竞争侧信道（RRMB）\cite{chiang2025reloadreload}。研究表明，在 SEV 体系结构中，不同地址空间标识符（ASID）的实体访问同一物理地址区域时会触发跨 ASID 的缓存刷新行为，从而形成可观测的时间差异。此外，当不同实体竞争访问同一内存块时也会产生明显的访问延迟。这些延迟差异可以被利用构造侧信道攻击，从而推断机密虚拟机的内存访问模式。

与此同时，Giner 等人在 Cohere+Reload 工作中深入剖析了这一现象的根源，指出 SEV-SNP 平台中用于维护宿主机密文视图与客户机明文视图一致性的硬件机制，是导致大范围（2KB）缓存被强制驱逐的根本原因。在 SEV-SNP 虚拟化中，虚拟机访问的是解密后的缓存数据，而虚拟机监控器在某些场景下可以访问同一物理地址对应的加密数据副本。为了保证两种数据视图之间的一致性，处理器会在访问其中一个副本时驱逐另一种副本，从而产生明显的访问时间差异，因此可以进一步通过测量访问延迟推断机密虚拟机对特定内存区域的访问行为。\cite{giner2025coherereload}

\subsection{国内外研究总结}
表\ref{tab:cvm_attacks}归纳了近年来针对机密计算虚拟机的攻击研究。现有研究主要集中在瞬态执行攻击、传统缓存侧信道攻击以及近年来逐渐受到关注的缓存一致性相关侧信道攻击等方向。研究表明，即使在机密计算环境中通过内存加密和硬件隔离机制保护了内存数据，处理器微架构中的共享资源仍可能泄露敏感信息。

\begin{table}[h]
  \centering
  \renewcommand{\arraystretch}{1.16}
  \caption{机密计算虚拟机环境中的典型攻击研究}
  \label{tab:cvm_attacks}
  \begin{tabular}{p{3.5cm}p{2.2cm}p{2.5cm}p{3.8cm}p{2cm}}
    \hline
    研究工作 & 平台 & 攻击类型 & 关键机制 & 粒度 \\
    \hline
    SEVered \cite{morbitzer2018severed}
    & SEV
    & 内存攻击
    & 页表重映射
    & 页面级 \\

    Spectre \cite{kocher2020spectre}
    & 通用 CPU
    & 瞬态执行攻击
    & 分支预测侧信道
    & 微架构级 \\

    Meltdown \cite{lipp2018meltdown}
    & 通用 CPU
    & 瞬态执行攻击
    & 权限检查绕过
    & 缓存行级 \\

    Foreshadow \cite{bulck2018foreshadow}
    & Intel SGX
    & 瞬态执行攻击
    & L1TF 缓存泄露
    & 缓存行级 \\

    CrossLine \cite{li2021crossline}
    & SEV
    & 内存隔离攻击
    & 地址空间转换绕过
    & 页面级 \\

    Prime+Probe \cite{osvik2006cache}
    & 云环境
    & 缓存侧信道
    & Cache set 竞争
    & Cache set \\

    Flush+Reload \cite{yarom2014flushreload}
    & 通用 CPU
    & 缓存侧信道
    & Flush+Reload
    & 缓存行级 \\

    Reload+Reload \cite{chiang2025reloadreload}
    & SEV-SNP
    & 缓存侧信道
    & 缓存一致性 / 内存竞争
    & 2KB / 64B \\

    Cohere+Reload \cite{giner2025coherereload}
    & SEV-SNP
    & 缓存一致性
    & Ciphertext/plaintext 冲突
    & 2KB / 256B \\

    RMPocalypse \cite{schlüter2025rmpocalypse}
    & SEV-SNP
    & 协议攻击
    & RMP特权协议缺陷
    & 页面级 \\

    Heracles \cite{schlüter2025heracles}
    & SEV-SNP
    & 密码算法攻击
    & 选择明文攻击
    & 页面级 \\

    WeSee / Heckler \cite{schluter2024wesee, schluter2024heckler}
    & SEV-SNP
    & 异常注入攻击
    & 恶意 \#VC 与中断
    & 指令级 \\

    CacheWarp \cite{zhang2024cachewarp}
    & SEV-ES/SNP
    & 故障注入攻击
    & 状态选择性回滚
    & 指令级 \\

    \hline
  \end{tabular}
\end{table}

总体来看，现有研究已经表明机密计算平台的硬件隔离机制仍可能被底层的缓存与内存共享行为的侧信道攻击打破。然而，针对此类新型侧信道的研究仍然存在明显不足。现有工作大多侧重于理想环境下的现象展示，且倾向于将机密计算环境下的高延迟信号单一归因为加密域间的缓存一致性冲突。

在真实高噪音机密虚拟机平台中，跨域一致性逐出信号极易被混淆，细粒度的缓存侧信道恢复面临巨大挑战。真实的攻击侧信道信号，往往是缓存状态失效与底层物理总线竞争叠加的复合结果。此外，现有的防御研究也大多针对传统缓存攻击，缺乏针对一致性与竞争复合机理的系统性权衡。因此，有必要围绕真实机密虚拟机环境中的复合缓存侧信道问题，从机制判别、可利用性评估以及缓解方法等方面开展更加深入的实证研究。

\section{研究内容与目的}
本文面向 AMD SEV-SNP 机密虚拟机环境，围绕底层缓存机制引起的缓存访问延迟侧信道展开研究。与传统针对共享缓存占用、共享内存页的攻击不同，是机密计算平台处理器在维持跨加密域数据一致性及处理底层资源竞争时所产生的微架构时间差异，以及这些时间差异在机密虚拟化场景下是否能够被攻击者稳定观测并进一步利用，形成稳定的侧信道攻击。本文希望从攻击机理、攻击实现、实验评估和防御分析四个层面，对缓存一致性侧信道在机密虚拟机中的安全影响进行系统研究。

\noindent 本文主要关注以下三个问题：
\begin{itemize}
  \item 在 AMD SEV-SNP 机密虚拟机环境中，缓存访问时间差异的原理，以及这些时间差异是否形成稳定的侧信道信号。
  \item 在系统噪音与底层硬件机制的干扰下，如何在真实环境中实现高精度的侧信道特征提取与攻击优化。
  \item 针对这类新型侧信道攻击，是否有可行的安全防御手段，以及如何在性能与安全之间取得平衡。
\end{itemize}

\noindent 围绕上述问题，本文将从以下三个方面展开研究：

\begin{enumerate}
  \item \textbf{研究复合缓存侧信道的形成机理并分析其攻击能力} \newline
    首先，本文将研究缓存侧信道的形成机理，并分析其可能具备的攻击能力。缓存侧信道的形成依赖于目标内存访问在微架构层面上产生的时间差异，而导致该差异的底层原因通常包括缓存一致性协议和物理内存竞争机制等多重因素。一方面，缓存一致性协议是现代多核处理器中保证不同核心间数据视图一致性的基础同步机制，当多个实体对同一内存区域进行访问时，硬件需要通过状态迁移、失效消息传播、缓存行回写与驱逐等操作维持一致性\cite{amd64manual}。这种一致性维护机制本身服务于系统性能和正确性，但在跨加密域访问时，其强制驱逐行为往往伴随着显著且可观测的时间惩罚。另一方面，内存竞争机制则是指当两个实体在同一时刻并发访问同一个物理连续内存块时，底层硬件（如内存总线或控制器）需要通过物理仲裁和资源调度来协调这种并发冲突，从而在物理链路上引入额外的竞争等待与访问延迟\cite{ge2018survey}。
    本文结合 AMD SEV-SNP 机密虚拟机的架构特征，深入剥离对比机密虚拟机平台下跨域一致性逐出与内存竞争的效应差异，系统性研究真实物理环境下导致访问高延迟的缓存侧信道信号来源及其复合叠加机理。通过架构分析和实验分析，本文将评估这些由复合机制引发的时间差异信号是否具有足够的稳定性和可区分性，从而在强噪声干扰下形成可实际利用的侧信道特征。在此理论与物理验证基础上，本文将进一步全面分析该类复合侧信道所具备的实际攻击能力，具体包括其在真实平台上的空间粒度限制、时间分辨率、信息泄露的适用对象，以及与传统 Prime+Probe、Flush+Reload 等纯缓存状态侧信道方法在攻击路径、底层先验条件上的本质异同。这些机理层面的界定与分析，将为后续面向真实场景的攻击工程化优化与防御机制探讨提供严谨的理论支撑。

  \item \textbf{研究复合缓存侧信道在机密虚拟机中的实际可利用性}\newline
    其次，本文将研究由上述多重底层机制引发的缓存侧信道在机密虚拟机中的实际可利用性。微架构状态变化在理论上可观测，并不意味着其在真实系统环境下具有稳定的攻击效果，因此需要通过系统性的实验验证其实际威胁。本文将在 AMD SEV-SNP 架构下构建相应的系统与攻击者模型，明确攻击者在权限、资源以及观测手段上的边界条件。在此基础上，设计面向真实机密虚拟机环境的完整攻击流程，涵盖目标地址定位、并发访问同步、时延测量与数据特征提取等关键步骤。真实物理部署环境中往往存在系统强噪声干扰，导致细粒度的侧信道信号恢复受限。针对这一现状，本文将通过实验评估该侧信道在不同系统负载与核心绑定策略下的稳定性，并通过信号降噪、特征提纯以及自适应分类等工程化优化机制提高其在正式场景下的可用性，有效规避底层硬件架构干扰与系统噪音引发的侧信道难以应用的问题。

  \item \textbf{研究针对复合缓存侧信道攻击的防御方法并评估其有效性}\newline
    最后，本文将研究针对该类新型缓存侧信道攻击的防御方法，并评估其实际缓解效果。针对传统的微架构侧信道攻击，现有的防御手段主要涵盖软件层面的时间扰动、硬件层面的资源隔离以及系统层面的访问控制等\cite{ge2018survey}。然而，对于由缓存一致性与物理内存竞争共同引发的新型侧信道，现有防御策略往往缺乏针对性的系统研究。为此，本文将在攻击机理与威胁边界探明的基础上，从硬件机制与系统软件两个层面探讨可行的防御思路。本文也将探讨不同防御方案在安全收益与系统性能开销之间的权衡，为机密计算平台的安全加固与架构演进提供参考依据。
\end{enumerate}

\section{创新点概述}
相较于已有研究主要侧重于特定攻击现象的发现与可行性验证，本文尝试从机理分析、攻击评估和防御讨论等方面，对 AMD SEV-SNP 机密虚拟机环境中的复合缓存侧信道问题进行较为系统的研究。本文的主要创新点概括如下：

\begin{enumerate}
  \item \textbf{在 AMD SEV-SNP 环境中对缓存侧信道的形成机理及其实验链路进行了验证与分析} \newline
    现有关于机密虚拟机缓存侧信道的研究，往往更关注攻击现象的展示与攻击效果的验证，而对相关侧信道在真实虚拟化环境中是否能够稳定形成、其底层实现条件是否成立缺乏系统分析。本文结合 AMD SEV-SNP 机密虚拟机的运行环境，在宿主机侧构建了不依赖虚拟机内部配合的验证模型，对宿主机与机密虚拟机之间由缓存一致性和内存竞争所引发的时序差异进行了验证，对比了机密虚拟机平台下跨域一致性逐出与物理内存竞争的效应差异，从体系结构和微架构角度分析了该类侧信道产生的原因及其影响因素。

  \item \textbf{提出了面向强噪声真实场景的复合缓存侧信道攻击优化方法} \newline
    传统高分辨率缓存侧信道方法通常建立在信号较强、环境较稳定的假设之上，假定攻击者可以实现细粒度的精准监控，但在真实机密虚拟机环境中，由于系统噪声、资源竞争以及平台结构特性的影响，细粒度信号恢复往往受到明显限制。针对这一问题，本文面向真实运行环境中的弱信号和高噪声特征，对复合缓存侧信道的攻击过程进行了优化，通过改进信号判别与攻击策略，提高了该类侧信道在复杂环境中的可利用性和稳定性。实验结果表明，经过优化后，该类攻击在真实机密虚拟机环境下仍能够取得较为稳定的识别效果，证明此类复合缓存侧信道具有进一步工程化利用的潜力。

  \item \textbf{结合攻击机制与实验结果，讨论面向复合缓存侧信道的防御思路} \newline
    现有针对机密计算平台的防御研究，主要集中于传统缓存侧信道、瞬态执行攻击以及一般性的资源隔离方法，对此类新型复合缓存侧信道问题缺乏专门分析。本文在机理分析和实验验证的基础上，从硬件机制、系统软件等角度，讨论了降低该类侧信道风险的可能方法，并进一步分析了相关防御措施在安全收益、性能开销和部署复杂度之间的权衡关系，为未来机密计算平台的安全加固与架构演进提供了参考依据。
\end{enumerate}

\section{论文结构}
本文共分为八个章节，各章节的主要内容如下：

第一章为绪论部分。本章首先介绍了本文的研究背景与意义，说明了云计算与机密计算快速发展背景下，AMD SEV-SNP 机密虚拟机在保护使用中数据安全方面的重要作用，并指出机密虚拟化环境仍然面临微架构侧信道攻击的潜在风险。随后，本文分析并总结了国内外与机密计算安全、微架构侧信道攻击，以及由缓存一致性等底层机制引发的新型侧信道相关的研究现状，分析了现有研究的主要进展与不足。在此基础上，引出了本文的研究问题、研究内容与研究目的，概述了本文工作的主要创新点，最后给出了全文的组织结构安排。

第二章为相关技术基础。本章主要介绍本文研究所依赖的相关背景知识。首先介绍机密计算的基本概念及其技术体系，重点分析 AMD SEV、SEV-ES 和 SEV-SNP 的架构演进与安全特性；随后介绍虚拟化环境中的内存管理机制，包括客户机物理地址到宿主机物理地址的映射关系、页表转换以及虚拟机监控器在内存管理中的作用；接着介绍多核处理器中的缓存层次结构、缓存一致性协议以及物理内存总线竞争机制，重点说明 MESI、MOESI 等协议的基本原理及状态转换过程；最后介绍微架构侧信道攻击的基本思想和典型方法，为后续章节分析该类缓存侧信道的形成机理与攻击实现奠定基础。

第三章为威胁模型与问题建模。本章首先给出本文所研究的攻击模型，明确 AMD SEV-SNP 机密虚拟机环境中的参与实体及其交互关系，包括受害虚拟机、攻击者实体、虚拟机监控器以及底层硬件平台。随后给出攻击者能力假设和安全边界，分析攻击者在权限、资源访问能力和观测手段上的限制。在此基础上，本章对侧信道时序信号的来源进行建模，研究缓存一致性状态变化、物理总线竞争与访存延迟之间的关联关系，并从理论上分析在机密虚拟机环境中构造该类侧信道攻击的可行性与约束条件。

第四章为真实环境下的侧信道机理辨析与攻击面分析。本章围绕缓存一致性与内存竞争机制所带来的微架构可观测行为展开分析。通过实验对比机密虚拟机环境下缓存一致性与物理内存竞争的效应差异，系统性研究真实物理环境下导致缓存访问延迟的信号来源及其叠加机理。随后，本章将该类侧信道与传统缓存攻击方法进行对比，分析其在攻击路径、信号粒度、先验条件和适用场景上的异同，从而明确该类攻击在机密计算环境中的特点。

第五章为攻击设计与实现。本章首先介绍本文实验平台和测试环境的搭建方式，包括 AMD SEV-SNP 机密虚拟机环境的配置、攻击者与受害者的部署方式以及实验所使用的关键软硬件组件。随后，本章给出基于复合缓存侧信道的攻击流程设计，包括目标地址定位、并发访问同步、时延测量与特征提取等关键步骤。在此基础上，详细介绍面向强噪声真实环境的攻击工程化优化细节，通过改进判别特征和攻击策略有效提高了真实物理环境下攻击成功率，并在不同实验环境下对攻击流程进行验证。最后，本章还将分析该攻击在不同实验条件下的稳定性、鲁棒性和适用范围，为后续实验评估提供实现基础。

第六章为实验评估与结果分析。本章主要对所提出攻击方法的有效性进行实验验证与量化分析。首先介绍实验评估所采用的指标，包括访问时延分布、攻击成功率、误判率以及不同负载条件下的稳定性表现。随后，本章从时序分布统计、不同系统负载影响、不同参数设置影响以及跨场景实验等多个角度，分析该类缓存侧信道在机密虚拟机环境中的实际可利用性，重点评估自动化选行策略与自适应动态阈值机制在真实强噪声环境下的特征提纯与准确率提升效果。此外，本章还将该方法与传统 Prime+Probe 等缓存攻击方法进行对比，分析其在攻击精度、可行性和适用条件上的差异，并进一步评估该类攻击对 SEV-SNP 安全模型所带来的实际影响。

第七章为防御机制与缓解策略。本章在前述攻击机理分析与实验结果的基础上，讨论针对该类新型缓存侧信道攻击的可能防御方法。首先分析现有微架构侧信道防御策略在机密虚拟机环境中的适用性与局限性；随后从硬件机制、系统软件两个层面，分别讨论降低一致性状态变化可观测性、通过隔离并发调度缓解物理总线争用，以及削弱攻击信号稳定性的缓解思路。在此基础上，本章进一步分析不同防御方案的性能开销和部署复杂度，探讨安全性与系统性能之间的权衡关系，为机密计算平台的安全加固提供参考。

第八章为总结与展望。本章对全文的研究内容与主要工作进行总结，归纳本文在缓存侧信道系统性分析、真实环境下的缓存攻击优化与实验验证，以及缓存侧信道防御方法探讨等方面取得的主要结论。结合本文研究过程中存在的限制与不足，分析还存在的改进空间，并对机密计算环境下微架构侧信道攻击与防御的未来研究方向进行展望。
